\documentclass[11pt,a4paper]{article}

\usepackage[utf8]{inputenc}

\usepackage[T1]{fontenc}

\usepackage[polish]{babel}

\usepackage{lmodern}

\usepackage{geometry}

\geometry{margin=27mm}

\usepackage[hidelinks]{hyperref}

\title{Etap edukacji a odroczony transfer samodzielnego osądu u uczniów korzystających z AI}

\author{}

\date{Szkic naukowy, wersja 0.1.1}

\begin{document}

\maketitle

\begin{abstract}

Tematem pracy jest Etap edukacji a odroczony transfer samodzielnego osądu u uczniów korzystających z AI. Tekst systematyzuje uzasadnienie, hipotezy i sposób weryfikacji zagadnienia. Proponowany wskaźnik: Modyfikacja efektu transferu przez etap/typ edukacji przy wspólnej skali wynikowej. Nie przedstawia się niezweryfikowanych wyników ani fikcyjnych cytowań.

\end{abstract}

\section{Wprowadzenie}

Samodzielność osądu i funkcje poznawcze zależą od jakości informacji oraz sposobu jej przedstawienia. Zagadnienie Etap edukacji a odroczony transfer samodzielnego osądu u uczniów korzystających z AI wymaga odróżnienia badanego konstruktu od jego operacjonalizacji i kontekstowych źródeł zmienności. Trzeba porównać możliwe wyjaśnienia i ustalić granice interpretacji dowodów.

\section{Cel i pytania badawcze}

Celem pracy jest rozstrzygnięcie zagadnienia: Etap edukacji a odroczony transfer samodzielnego osądu u uczniów korzystających z AI. Proponowany wskaźnik: Modyfikacja efektu transferu przez etap/typ edukacji przy wspólnej skali wynikowej. Hipotezy i ich kierunki muszą zostać określone przed analizą danych lub krytyczną syntezą źródeł.

\section{Materiał i metody}

Podłużne zadania osądu bez AI zagnieżdżone w klasach różnych etapów kształcenia. Pierwszorzędowy wskaźnik: Jakość odroczonego transferu samodzielnego osądu. Planowany projekt nie obejmuje jeszcze rekrutacji ani pomiarów. Przed ich rozpoczęciem należy odrębnie zatwierdzić kryteria doboru, liczebność próby, etykę, ochronę danych i prawa do materiałów.

\section{Plan analizy}

Przewidziano model wielopoziomowy z korektą poziomu bazowego i efektu klasy. Zamrożenie estymandu, kryteriów jakości i analiz wrażliwości nastąpi przed dostępem do danych; na tym etapie nie ma oszacowań ani wyników.

\section{Wyniki}

W niniejszym pliku nie umieszczono wyników empirycznych, statystyk ani nieprzeprowadzonych analiz. Sekcja wymaga uzupełnienia po uzyskaniu i weryfikacji materiału dowodowego.

\section{Dyskusja i ograniczenia}

Ograniczenia obejmują możliwe błędy pomiaru, dobór próby, alternatywne mechanizmy i ograniczoną możliwość uogólniania. Korelacja nie dowodzi przyczynowości; poprawność techniczna nie jest dowodem znaczenia klinicznego.

\section{Wnioski}

Zdefiniowano ramy opracowania, ale nie sformułowano wniosków wykraczających poza dostępny stan dowodów.

\section*{Bibliografia}

Weryfikacja literatury i przypisanie cytowań do stwierdzeń są kolejnym etapem prac.

\end{document}

